\documentclass[aspectratio=169,12pt]{beamer}
\usepackage[T1]{fontenc}\usepackage[utf8]{inputenc}\usepackage[italian]{babel}
\usepackage{amsmath,graphicx,booktabs,tikz}\usetikzlibrary{arrows.meta,positioning,calc}
\definecolor{deepblue}{RGB}{20,52,92}\setbeamercolor{structure}{fg=deepblue}
\setbeamercolor{frametitle}{fg=white,bg=deepblue}\setbeamertemplate{navigation symbols}{}
\setbeamertemplate{footline}{\hfill\insertframenumber/\inserttotalframenumber\hspace{1em}\vskip.5em}
\hypersetup{pdfauthor={Alessio de Dato; Valerio Rebechi},pdftitle={SEA: incertezze parametriche e concentrate}}
\newcommand{\diagram}[1]{\centering\resizebox{.86\linewidth}{!}{\input{../schemi/#1}}\par}
\newcommand{\plotdata}[1]{\centering\includegraphics[height=.64\textheight,width=.94\linewidth,keepaspectratio]{../../figure/#1}\par}
\title{Manipolatore SEA a due giunti}\subtitle{Due rappresentazioni dell'incertezza degli attuatori}
\author{Alessio de Dato -- 635256\\Valerio Rebechi -- 722224}\date{25 ottobre 2026}
\begin{document}\begin{frame}[plain]\titlepage\end{frame}
\begin{frame}{Sistema SEA ed equilibrio}
\diagram{sea}
\small $\bar q=[\pi/2,0]^T$, 8 stati meccanici + 6 drive. Encoder link/rotori per LQG; q per i robusti.
\end{frame}
\begin{frame}{Attuatore nominale}
\diagram{nominale}
\[G_{a0}=\frac{40^2}{s^2+64s+1600}\frac{1-.015s/2}{1+.015s/2}\]
\small Banda 34.78 rad/s; zero Padé a +133.33 s$^{-1}$. Nessuna duplicazione di Jm/Bm.
\end{frame}
\begin{frame}{Scelta di pochi parametri incerti}
\small
\begin{center}\begin{tabular}{lrr}\toprule
Parametro&Intervallo&Scost. max a 40 rad/s\\\midrule
$\omega_n$&36--44 rad/s&16.41\%\\$\zeta$&.72--.88&11.11\%\\$\tau$&12--18 ms&11.18\%\\\bottomrule\end{tabular}\end{center}
Cinque valori per parametro, uno alla volta. S1 approvato: wn/zeta incerti,
tau noto a 15 ms; meccanica nominale nel confronto centrale. M4 conservato separatamente.
\end{frame}
\begin{frame}{A parametrico e LFT manuale}
\diagram{parametrico}
\[\Delta_A=\operatorname{diag}(\delta_{\omega1}I_2,\delta_{\zeta1},\delta_{\omega2}I_2,\delta_{\zeta2})\]
\small Quattro parametri reali indipendenti; sei porte. Le ripetizioni di wn sono correlate.
\end{frame}
\begin{frame}{B concentrato diagonale}
\diagram{concentrato}
\[W_I=.6165601\frac{s}{s+36},\quad \|\Delta_i\|_\infty\le1\]
\small Due dinamiche SISO indipendenti. Il blocco pieno aggiungerebbe accoppiamenti non previsti.
\end{frame}
\begin{frame}{Copertura e conservativita}
\plotdata{01_attuatori.png}
\small Prova analitica sul box: $\mathcal A\subseteq\mathcal B$. Massimo campionato $\|E/W_I\|_\infty=.559738$.
L'inclusione unitaria non ordina automaticamente tutte le espansioni dei margini.
\end{frame}
\begin{frame}{Segnali e pianta generalizzata}
\diagram{generalizzata}
\small Stessi W1, W2 e Mo per A/B; d e una coppia sui link. Norma pesata e target temporali restano distinti.
\end{frame}
\begin{frame}{LQG e integratore misurato}
\diagram{lqgi}
\small LQG: $u=-K\hat x+N_rr$. LQGI: $\dot\xi=r-q_{mis}$, $u=-K_x\hat x-K_i\xi$.
Stesso Kalman; Q/R e covarianze consolidate. LQG Ts=.675 s; LQGI Ts=5.05 s con reiezione DC del disturbo.
\end{frame}
\begin{frame}{H-infinity: metodi e problemi}
\begin{center}\small\begin{tabular}{lrr}\toprule
Metodo & Ordine & Norma proprio problema\\\midrule
mixsyn & 18 & 0.959295\\
hinfsyn stesso P1 & 18 & 0.955559\\
hinfsyn P2 & 22 & 0.0475849\\
hinfstruct PID P2 & 4 & 0.800896\\
\bottomrule\end{tabular}
\end{center}
\small Le prime due righe risolvono lo stesso P1.
Il 2-DoF e il PID usano P2: gamma di problemi diversi non costituisce una classifica.
\end{frame}
\begin{frame}{NS, NP, RS e RP}
\diagram{robusta}
\small NS: nominale internamente stabile. NP: norma comune sotto uno.
RS: stabilita sull'intera famiglia. RP: anche prestazione.
Un risultato campionario non certifica RS/RP.
\end{frame}
\begin{frame}{Stessi controllori, rappresentazione diversa}
\begin{center}\small\begin{tabular}{lrr}\toprule
K & RS lower A & RS lower B\\\midrule
LQG & 1.728 & 0.9914\\
LQGI & 1.68 & 0.957\\
Hinf1 & 2.586 & 1.046\\
Hinf2 & 0.5655 & 0.0332\\
PID & 0.2945 & 0.1252\\
Kmu & 3.353 & 1.898\\
\bottomrule\end{tabular}
\end{center}
\small LQG/LQGI: PASS su A, FAIL su B. Hinf2/PID fragili in entrambi.
Kmu conserva RS/RP. Il concentrato copre un insieme piu ampio del box fisico.
\end{frame}
\begin{frame}{Due D-K dedicate, stessi pesi}
\plotdata{04_robustezza.png}
\small Quattro iterazioni massime e stesso fit [2,1]. Il K-step non sostituisce la post-analisi.
\end{frame}
\begin{frame}{Sintesi e verifica incrociata}
\begin{center}\small\begin{tabular}{lrrr}\toprule
K & Ordine & RP lower A & RP lower B\\\midrule
Kmu & 78 & 1 & 1.448\\
Mu\_A & 88 & 3.873 & 2.03\\
Mu\_B & 32 & 4.226 & 2.625\\
\bottomrule\end{tabular}
\end{center}
\small Kmu su A: bound da inclusione, margine diretto inconcludente.
Mu\_B ha ordine 32 contro 88 di Mu\_A; nessuna riduzione adottata.
\end{frame}
\begin{frame}{Dieci impianti fisici comuni}
\plotdata{orale_impianti.png}
\small Hinf2 stabile 3/10, PID 7/10; gli altri 10/10. Stessi campioni per tutti. Solo le risposte stabili sono tracciate.
\end{frame}
\begin{frame}{Sezione empirica di attrazione}
\plotdata{orale_ras.png}
\small Quindici punti per K; Delta theta=Delta q, altri stati nulli. Verde convergente, rosso non convergente entro T, grigio inconcludente.
Mu\_A: 6/15; Mu\_B: 15/15. Nessun bordo certificato.
\end{frame}
\begin{frame}{Rumore e risultati conservati}
\small I nuovi K non penalizzano esplicitamente gli encoder:
\[\|T_{uv}\|_\infty:\quad K_\mu\simeq2.94\,10^4,\quad K_{\mu B}\simeq1.36\,10^6.\]
La formulazione storica con rumore produceva MuN1 a 82 stati: RS/RP su M4,
picco Simulink non lineare 108.526 $\to$ 2.251 N m, ma offset q1 peggiore.
\medskip\par Famiglie, pesi ed evidenze restano distinti. I nuovi SLX saranno forniti dagli studenti.
\end{frame}
\begin{frame}{Conclusioni del confronto A/B}
\begin{itemize}
\item A conserva significato fisico e correlazioni; B aggiunge dinamiche ammesse e conservativita.
\item Su K fissi, B riduce qui diverse garanzie. Nella sintesi finita, B produce il candidato piu compatto.
\item Mu\_B supera RS/RP su A/B; Ts=1.985 s e DC=5.454\% non rispettano tutti i target.
\item Rumore, ordine e comportamento non lineare impediscono una scelta basata solo sui margini.
\end{itemize}
\end{frame}
\end{document}
